\section{Other targets and categorical representability}
\label{sec:targets}

Write $\mathcal C_f$ for a skeleton of finite categories, with functors
identified when they are joined by a finite zigzag of natural transformations.
Thus $h_K(C)=\pi_0\Fun(C,K)$, whereas $F_Y(C)=[BC,Y]$.
These are different expressions even when $K$ is finite.
A \emph{Postnikov piece} here is a connected homotopy type with
$\pi_i=0$ for all sufficiently large $i$; the remaining homotopy groups
need not be finite.  A disconnected space is \emph{componentwise a
Postnikov piece} if this holds in each component, without a common
truncation bound.

\subsection{Audit of the representation criterion used below}

The following audits concern the specified arguments in CATHOT D.3--D.5
and INDLOCAL B.8--B.9.  They are included to specify exactly which
comparisons are being used.

\begin{proposition}[Verified form of R1 D.3--D.4]
\label{CD:sequential}
For every homotopy type $Y$, the presheaf $F_Y$ on $\mathcal C_f$ is a
filtered colimit of finite-category representables.  The following
conditions are equivalent:
\begin{enumerate}[label=(\roman*)]
\item $F_Y$ is a colimit of a sequence of such representables;
\item $[BC,Y]$ is countable for every finite category $C$;
\item $F_Y$ is represented on finite tests by a countable locally finite
category $K$, where ``locally finite'' means that every finitely generated
subcategory has finitely many arrows.
\end{enumerate}
\end{proposition}
\begin{proof}
We check the steps of the supplied proof.  An object of its element
category is $(C,c)$ with $c\in[BC,Y]$, and an arrow to $(D,d)$ is a
categorical homotopy class $u:C\to D$ with $c=dBu$.
The empty category supplies an object; disjoint union supplies an upper
bound for two objects.  For parallel arrows represented by $u,v:A\to B$,
form the finite category with objects those of $A$ and $B$ and with two
tagged copies of $B(u(a),b)$ and $B(v(a),b)$ as cross arrows $a\to b$.
The tagged identities give natural transformations joining the two
composites of $u,v$ with $B\to D$.  Its realization is the double mapping
cylinder of $Bu,Bv$: the categorical mapping-cylinder comparison is
obtained from the same tagged arrows and is a weak equivalence on each
half, relative to the copies of $BA$ and $BB$.  A homotopy between
$bBu$ and $bBv$ therefore extends $b:BB\to Y$ over $BD$.
This proves the equalization axiom, with the stated arrow $B\to D$.

The density map sends a functor $f:T\to C$ indexed by $(C,c)$ to
$cBf$.  Its surjectivity is witnessed by $\id_T$ at $(T,t)$.
If two representatives have the same image $t$, the arrows from
$(T,t)$ to their indexing objects identify both with this same identity
representative.  Thus density gives precisely $F_Y$, not a different
presheaf on a localization.

There are countably many finite composition tables and finitely many
functors between two fixed finite categories.  Under (ii) the element
category consequently has countably many objects and arrows.  The usual
cofinal-sequence argument is valid here: enumerate its objects and
arrows; successively take a common successor and equalize the next
finite set of parallel-arrow requirements.  A diagonal enumeration of
the requirements ensures that every comma category from an object to
the resulting sequence is nonempty and connected.  This is a cofinal
functor from $\mathbb N$, so it gives (i).  Conversely, each $h_{K_n}(C)$
is finite when $C,K_n$ are finite, proving (i)$\Rightarrow$(ii).

Choose actual functors for consecutive transition morphisms; composing
them defines an actual sequence since this index has no extra relations.
A functor from finite $C$ into $K=\colim_nK_n$, a natural transformation,
and an equality between either sort of finite data all occur at a common
finite stage.  Hence
\[
 \Fun(C,K)=\colim_n\Fun(C,K_n),\qquad
 h_K(C)=\colim_n h_{K_n}(C).
\]
The second equality uses only that a zigzag is finite.  The colimit $K$
is countable.  A finite generating set of its arrows lifts to one stage;
move farther so that every equality between their finitely many endpoint
objects already holds.  Every composable word in the generators now
lifts to that finite stage, so the generated subcategory is finite.
This verifies the endpoint-identification step in D.4 and proves (iii).
Conversely enumerate the arrows of a countable locally finite category
and use the finite subcategories they generate.  Their union and the
same finite-data argument give (i).
\end{proof}

\begin{proposition}[Verified finite bound in R5 B.9]
\label{CD:uniformfinite}
A filtered ind-object whose stages have at most $m$ arrows, for a fixed
finite $m$, is isomorphic in $\Ind(\mathrm{Cat}_{\mathrm{hot}})$ to a
constant finite category with at most $m$ arrows.
\end{proposition}
\begin{proof}
Here are the checks of the supplied proof.  The full subcategory
$\mathcal D$ containing one representative of each category of size at
most $m$ has finitely many objects and arrows.  If
$F=\colim_i h_{K_i}$ on $\mathcal D$, then
$|F(A)|\le\max_T|\mathcal D(A,T)|$: any larger finite collection would
have distinct representatives at one common stage.  Its finite element
category is filtered because finite data and equalities lift to stages.
A cocone on its identity diagram expresses $F$ as the image of an
idempotent in $\mathcal D(T,T)$.  The full inclusion of $\mathcal D$
induces a fully faithful inclusion of ind-categories by the
limit--colimit formula for morphisms, so this splitting concerns the
original ind-object as well.

Represent the idempotent class by an endofunctor $f:T\to T$.
The finite monoid of endofunctors has eventually periodic powers.
Choose a sufficiently large positive multiple $r$ of the period;
$f^{2r}=f^r$ strictly, while $[f^r]=[f]$.  The objects and arrows fixed
by $f^r$ form a subcategory $L$ with $sj=\id_L$ and $js=f^r$ for its
inclusion $j:L\to T$ and restriction $s:T\to L$.
This supplies exactly the required splitting and size bound.  The case
$m=0$ consists only of the empty category.
\end{proof}

For completeness, the finite-test part of R5 B.8 has the following
direct check.  A filtered colimit of representables takes the empty
category to a point and disjoint unions to products.  An equality of
two restrictions of an element to $A$ is witnessed at one stage by a
finite zigzag of transformations.  Insert its finite fence in the
double-mapping-cylinder category above.  The resulting finite category
maps by a weak equivalence to the unsubdivided cylinder, relative to
$B$; an invariant presheaf transfers the extension across this map.
Conversely those product and extension conditions are exactly the
nonemptiness, upper-bound, and equalization conditions for the element
category.  Density then gives the presentation by finite stages.
This verifies B.8 at $\kappa=\aleph_0$ without treating a finite nerve
as a finite simplicial set.

\begin{lemma}[The cardinal-controlled acyclic comparison used in R5 B.7]
\label{CD:acycliccomparison}
Let $W$ be connected with a simplicial model of infinite cardinality at
most $\nu$, and let $\mu$ be infinite.  There is an acyclic map
$t:BG\to W$, with a surjective fundamental-group map $p:G\to\pi_1W$,
such that
\[
 |G|\le\max(\nu,\mu^+),\qquad
 \Hom(G,L)=\Hom(\pi_1W,L)\quad(|L|\le\mu),
\]
where the equality means the bijection induced by $p$.
If $Y$ has fundamental groups of cardinality at most $\mu$, then
$t^*:[W,Y]\to[BG,Y]$ is bijective.
No bound on the dimension of $BG$ is asserted.
\end{lemma}
\begin{proof}
Here we verify the original ingredients INDLOCAL A.1--A.3 used by
B.7.  Put $\lambda=\mu^+$.  Every group $L$ of size at most $\lambda$
embeds in a simple acyclic group of size $\lambda$, by the following
two explicit constructions.  For simplicity start with $L*F_\lambda$.
For each $g\ne1,h$ in a current stage, take a new free letter $t$ and
\[
 w=gtg^{-1}t^{-1}gt^2g^{-1}t^{-2},\qquad s^{-1}ws=wh.
\]
Both $w$ and $wh$ have infinite order: their cyclically reduced
free-product lengths are respectively eight, and eight or six.
The HNN normal form preserves the original group, and the equation
puts $h$ in the normal closure of $g$.  Repeating for all pairs and
then countably many times yields a simple group of size $\lambda$.

For acyclicity embed $L$ as the first axis of $L\times L$, write the
axes as $j,k$, and adjoin stable letters conjugating $j(l)$ respectively
to $j(l)k(l)$ and to $k(l)$.  These are HNN extensions over embedded
copies of $L$.  Suppose a map into $L$ kills positive homology below
degree $n$ over a field.  In degree $n$ its image under the diagonal
has only the two axis terms, by K\"unneth.  The two axes and the
diagonal become conjugate, so this image satisfies $a=2a$ and is zero.
Induction shows that successive such embeddings kill every positive
homology degree in their countable union.  Vanishing over $\mathbb Q$
and over all $\mathbb F_p$, with the universal coefficient sequence
in the next degree, gives integral acyclicity.  Alternating the simple
and acyclic envelopes yields both properties; their cofinal stages
verify simplicity and homology separately.  All cardinalities remain
$\lambda$.

Replace $W$ by a simplicial complex $K$ of size at most $\nu$.
The original construction assigns a space $T(L)$ to every subcomplex
$L\subseteq K$, compatibly with inclusions.  On graphs use $T(L)=|L|$.
Inductively, for each $n$-simplex $\sigma$, $n\ge2$, embed
$\pi_1T(\partial\sigma)$ in a simple acyclic group $C_\sigma$ of
size $\lambda$.  Its group has size at most $\lambda$, since the
boundary has finitely many faces.  Attach the mapping cylinder of
$T(\partial\sigma)\to BC_\sigma$ along its free boundary end.
Send the cylinder radially into $\sigma$ and $BC_\sigma$ to its
barycentre.  This defines actual compatible maps $T(L)\to|L|$.

The asphericity and injection induction is checked as follows.
Every attachment is a graph of aspherical spaces with injective edge
groups.  Normal forms inject the vertex groups; the universal cover
is a tree of contractible spaces, and its finite subtrees are
contractible.  Compactness of spheres gives contractibility of the
union.  For the subcomplex assertion $L\subseteq K$, first form
$T(K^{(n-1)})\cup_{T(L^{(n-1)})}T(L)$; both intersection maps inject
on fundamental groups.  Attach the remaining $n$-simplex cylinders
afterwards.  Their boundary groups already inject into the lower
skeleton, so all needed injections persist.  This proves the stronger
subcomplex induction rather than merely asphericity of the full space.

For a local system on $|L|$, its pullback is constant over every new
simplex.  The relative homology sequences for a cylinder and for
$(\sigma,\partial\sigma)$ agree: the new vertex space is acyclic,
and the boundary comparison is an isomorphism by induction.  Excision
and the skeletal sequence prove the homology comparison with every
local system, also in the union over dimensions.  Let $G=\pi_1T(K)$.
The kernel of $p:G\to\pi_1|K|$ is exactly the normal closure of the
$C_\sigma$.  In dimension two their removal imposes precisely the
triangular boundary relators; in higher dimensions the boundary group
modulo its previous cone groups is already
$\pi_1\partial\sigma=1$.  This proves the kernel assertion and
surjectivity.  Simplicity and $|C_\sigma|>\mu$ make every map to a
group of size at most $\mu$ kill each $C_\sigma$, hence the kernel.
The spaces and groups have size at most $\max(\lambda,\nu)$.

Finally we check the mapping assertion of A.2.  For any acyclic
$t:X\to W$ and $f:X\to Y$ killing $\ker\pi_1t$, form
$D=Y\cup_X^hW$ in the component containing $f(X)$.
Van Kampen gives $\pi_1Y\cong\pi_1D$.  Local-coefficient
Mayer--Vietoris gives a homology equivalence on universal covers;
the homological Whitehead theorem makes $Y\to D$ an equivalence.
A homotopy inverse extends $f$ across $t$.  For injectivity, the fold
$W\cup_X^hW\to W$ is likewise a weak equivalence, by the same
fundamental-group and local-homology calculation.  The two inclusions
of $W$ are equal in the homotopy category; hence equal restrictions
imply equal maps from $W$.  In the present construction every map
$BG\to Y$ kills the kernel, so both assertions apply.  This verifies
the precise comparison used in B.7 independently of any
two-dimensional-resolution assertion.
\end{proof}

\begin{proposition}[The sharp infinite threshold in R5 B.7]
\label{CD:boundedtests}
Let $\mu,\kappa$ be infinite cardinals.  For a small filtered
ind-category $K=(K_i)$ with $|\Mor K_i|\le\mu$, invariance on all
tests with fewer than $\kappa$ arrows forces the canonical global
ind-isomorphism $K\to(\Pi_1K_i)$ for every such $K$ exactly when
$\mu^+<\kappa$.
For $\kappa\le\mu^+$, every small homotopy type has an invariant
presentation on these tests with stages of at most $\mu$ arrows;
$K(\Z/2,2)$ gives one with no global ind-groupoid isomorphism.
\end{proposition}
\begin{proof}
We check the bounded localization in the supplied proof.  For infinite
$\nu$, nerves and categories of simplices preserve the bound $\nu$.
The last-vertex comparisons give their localization equivalence with
types having simplicial models of at most $\nu$ simplices.  A Kan
replacement attaches at most $\nu$ finite horns in countably many
rounds.  Maps and homotopies are therefore represented within the same
bound.  This proves full faithfulness as well as essential surjectivity.

Put $\nu=\mu^+<\kappa$.  On this bounded category let
$H=\colim_i h_{K_i}$ and let $F$ be its descent to homotopy types.
The restricted Yoneda calculation is
\[
 \operatorname{Nat}(\colim_i h_{BK_i},T)
 =\lim_iT(BK_i)
 =\operatorname{Nat}(H,q^*T)
 =\operatorname{Nat}(F,T),
\]
where $q$ is the localization.  Thus
$F(X)=\colim_i[X,BK_i]$, with a functor sent to its classifying map.
All stages lie in the bounded category, which is essential to this
calculation.

Apply Lemma~\ref{CD:acycliccomparison} to each component of $BK_i$.
Its resolution $BG\to BK_i$, and the groupoid $B\pi_1BK_i$, stay
within the bound $\nu$.  For every $j$, precomposition with its
quotient $G\to\pi_1BK_i$ gives an isomorphism of functor categories
into $K_j$: an action has image in an automorphism group of size at
most $\mu$, so factors uniquely; surjectivity of the quotient makes
the two natural-transformation intertwining conditions identical.
The acyclic comparison also makes $[BK_i,BK_j]\to[BG,BK_j]$
bijective.  The triangle formed by $t$, the quotient, and the unit
$BK_i\to B\Pi_1K_i$ commutes in homotopy, because its target is a
groupoid nerve.  Taking colimits proves
$H(\Pi_1K_i)\to H(K_i)$ bijective.

Take the limit over $i$ in the ind-Hom formula.  It supplies
$r:(\Pi_1K_i)\to K$ with $r\eta=\id_K$.  Every groupoid, and hence
every filtered ind-groupoid, is orthogonal to all units
$X\to\Pi_1X$: functors and transformations extend uniquely over
the inverted arrows.  Thus precomposition with $\eta$ is also
injective for maps to $(\Pi_1K_i)$, giving $\eta r=\id$.
This is the canonical global isomorphism claimed in B.7.

If $\kappa\le\mu^+$, repeat the element-category proof of
Proposition~\ref{CD:sequential} on the skeleton of categories with
at most $\mu$ arrows.  Finite coproducts and the tagged double
mapping cylinder preserve this bound, so the category is small and
filtered and density gives $[B-,Y]$ on all these tests.  Every test
of size less than $\kappa$ is included.  For $Y=K(\Z/2,2)$ a finite
face poset of $S^2$ gives two classes, whereas a point gives one.
Every ind-groupoid takes the map from that face poset to a point
to a bijection.  This rules out a global ind-groupoid isomorphism
and includes the equality case $\kappa=\mu^+$.
\end{proof}

The same localization size check proves the uncountable part of R5 B.8:
for $\kappa>\aleph_0$, each individual model of size less than
$\kappa$ can be replaced within
$\max(\aleph_0,\lambda)<\kappa$.  The element-category conditions
above are then exactly filteredness on $\mathrm{Hot}_{<\kappa}$.
Their invariant restrictions are its small ind-objects.  No regularity
of $\kappa$ is required, since only finite constructions and the
replacement of one bounded model at a time are used.

\subsection{The complete countability criterion}

\begin{lemma}[The wedge detects every unbounded Postnikov tower]
\label{CD:wedgecriterion}
Let $W=\bigvee_{j\ge2}S^j$, with its CW wedge topology, and let $Y$
be connected.  If $\pi_j(Y)\ne0$ for infinitely many $j$, then
$|[W,Y]|\ge2^{\aleph_0}$.  This assertion makes no countability
assumption on $\pi_1Y$.
\end{lemma}
\begin{proof}
The coproduct property of the pointed CW wedge gives
$[W,Y]_*\cong\prod_{j\ge2}\pi_j(Y)$, and the passage to free classes
is the orbit quotient by the usual simultaneous action of $\pi_1Y$.
Choose a nonzero $a_j$ for each index in an infinite set $J$ of
nonvanishing groups.  For every $T\subseteq J$, take the tuple with
coordinate $a_j$ at $j\in T$ and zero elsewhere.  An automorphism of
a homotopy group fixes zero and cannot carry a nonzero element to zero.
Consequently different supports $T$ lie in different orbits.  The
displayed construction is an injection from $\mathcal P(J)$ to free
homotopy classes.
\end{proof}

\begin{theorem}[Countability for arbitrary homotopy types]
\label{CD:alltargets}
The following are equivalent for a CW homotopy type $Y$.
\begin{enumerate}[label=(\roman*)]
\item $[BC,Y]$ is countable for every finite category $C$.
\item $\pi_0Y$ is countable, every component $Y_c$ is a Postnikov piece,
and $[A,Y_c]$ is countable for every finite connected CW complex $A$.
If $Y_c$ is an $n_c$-type, only $\dim A\le n_c+1$ is necessary here.
\item There are a countable CW complex $Z$, componentwise a Postnikov
piece, and a map $q:Z\to Y$ that induces bijections
$[A,Z]\to[A,Y]$ for every finite CW complex $A$.
\end{enumerate}
Under (iii), $q$ also induces bijections on all finite-category tests.
Together with Proposition~\ref{CD:sequential}, this characterizes exactly
the types represented on finite tests by one countable locally finite
category.
\end{theorem}
\begin{proof}
The six-element monoid of Part~A has classifying space $W$.
Thus (i) and Lemma~\ref{CD:wedgecriterion} force every component to be
truncated.  A point detects countability of the component set.
Every finite CW complex is homotopy equivalent to a finite simplicial
complex, whose face poset is a finite category; this proves the remaining
part of (ii).

If $Y_c$ is an $n_c$-type, restriction induces a bijection
\[
 [X,Y_c]\longrightarrow[X^{(n_c+1)},Y_c]
 \tag{\ref{CD:alltargets}.1}
\]
for every CW complex $X$.  Indeed, all subsequent attaching-sphere
obstructions lie in $\pi_kY_c$ for $k>n_c$, and the obstructions to
extending a homotopy relative to its two ends lie in the same vanishing
range.  The cell-by-cell extensions over all dimensions are continuous
by the weak topology.  For a finite category $C$, its nerve has only
finitely many cells in each degree.  The skeleton in this display is
therefore a finite CW complex.  Since $BC$ has finitely many components,
(ii) implies (i), including the absence of a uniform bound on the $n_c$.

To prove (ii)$\Rightarrow$(iii), choose a CW model for $Y$.
There are countably many finite simplicial complexes up to isomorphism.
Choose a representative map for every element of every $[A,Y]$ for
these complexes.  Their images lie in finite subcomplexes, so they are
contained in a countable subcomplex $Z_0$ of $Y$.
Given a countable subcomplex $Z_r$, each $[A,Z_r]$ is countable:
finite-domain simplicial approximation expresses it as a countable
union of finite or countable combinatorial map sets.
For every pair of such classes that become equal in $Y$, choose one
homotopy and enlarge to a countable subcomplex $Z_{r+1}$ containing its
image.  Include the paths required for $A$ a point as well.
Set $Z=\bigcup_{r\ge0}Z_r$.  Maps and homotopies with finite CW domain
have image in finite subcomplexes.  Thus the construction makes
$[A,Z]\to[A,Y]$ both onto and one-to-one for every finite $A$.
The component sets agree.  If $Y_c$ is an $n_c$-type, injectivity for
$S^j$, $j>n_c$, implies that every element of $\pi_j(Z_c)$ lies in the
orbit of zero and hence is zero.  Thus $Z$ has the required truncations.

Conversely (iii) first implies countability on finite CW tests.
It also implies that the corresponding component of $Y$ has exactly
the vanishing range of the homotopy groups of $Z$: a singleton
$[S^j,Y_c]$ forces $\pi_jY_c=0$, since the orbit of zero is a singleton.
Equation~(\ref{CD:alltargets}.1), on both sides of $q$, then gives the
bijection on finite-category tests and (i).
\end{proof}

\begin{corollary}[Targets with countable homotopy groups]
\label{CD:countabletargets}
For a countable CW homotopy type $Y$, the criterion is simply:
each component is a Postnikov piece.  In particular this is the exact
criterion for any finite CW complex $Y$.
More generally, if all component fundamental groups are countable,
the criterion is that there are countably many components, all higher
homotopy groups are countable, and each component is truncated.
\end{corollary}
\begin{proof}
Maps from a finite CW complex to a countable CW complex have countably
many homotopy classes, as used in the preceding proof.  Alternatively
one can count the finitely many obstruction stages directly: the
fundamental group of $BC$ is finitely generated; each local-coefficient
cochain group in a fixed degree is a finite product of countable groups.
For the more general assertion, the countability of $[S^j,Y_c]$ and
of $\pi_1Y_c$ implies countability of $\pi_jY_c$, since each orbit
is countable.  A type with countably many components and countable
homotopy groups admits a countable CW model, constructed by attaching
countably many cells at each successive homotopy stage.
\end{proof}

\begin{remark}[Why the qualification on fundamental groups matters]
\label{CD:largegroups}
Let $\Lambda$ be any infinite set, and let $A_{\mathrm{fin}}(\Lambda)$
be its finitary alternating group.  For a countably infinite
$\Lambda_0\subseteq\Lambda$, inclusion gives bijections
\[
 [BC,B A_{\mathrm{fin}}(\Lambda_0)]
 \longrightarrow [BC,B A_{\mathrm{fin}}(\Lambda)]
\]
for every finite $C$.  Indeed $\pi_1BC$ is finitely generated, so every
image acts on a finite support and is conjugate into $\Lambda_0$.
A conjugacy between two such finite images can be reproduced on the
union of their supports inside $\Lambda_0$; two spare letters correct
the parity.  This proves injectivity as well as surjectivity.
For uncountable $\Lambda$, the target fundamental group is uncountable.
Thus replacing Theorem~\ref{CD:alltargets} by a requirement of countable
homotopy groups would be false.
\end{remark}

\begin{corollary}[The targets requested in C1]
\label{CD:targetlist}
For every finite category $C$:
\begin{enumerate}[label=(\roman*)]
\item $[BC,S^1]$, $[BC,K(\Z,n)]$, and $[BC,K(\Z/p,n)]$ are
countable for $n\ge1$ and primes $p$; the last set is finite.
\item $[BC,BG]$ is countable for every finite or countable group $G$;
it is finite when $G$ is finite.
\item For every $n\ge2$ there is a six-arrow monoid $M$ with
$|[BM,S^n]|=2^{\aleph_0}$.  The same $M$ has
$|[BM,\mathbb CP^2]|=2^{\aleph_0}$.
\item A finite CW target has countable values on all finite-category
tests exactly when each of its components is a Postnikov piece.
\end{enumerate}
Each positive countability assertion is equivalent to the existence of
a sequence of finite categories and of a single countable locally finite
representing category.  Each negative assertion rules out both.
\end{corollary}
\begin{proof}
For an abelian group $A$, $[BC,K(A,n)]=H^n(BC;A)$.
The normalized cochain group in degree $n$ is a finite product of $A$.
This proves (i), including R1 D.5 by its stated finite-cochain argument.
Maps to $BG$ are homomorphisms $\pi_1BC\to G$ modulo conjugacy,
component by component; a finite set of generators gives (ii).
Serre's theorem gives infinitely many nonzero homotopy groups of every
$S^n$, $n\ge2$ \cite[\S3, Th\'eor\`eme 10 and its corollary,
pp.~217--219]{CDSerre1953}.  The bundle
$S^1\to S^5\to\mathbb CP^2$ identifies
$\pi_j\mathbb CP^2$ with $\pi_jS^5$ for $j\ge3$
\cite[Example~4.45]{CDHatcher2002}.  Lemma~\ref{CD:wedgecriterion}
and Part~A give the lower bounds in (iii).  The groups in question are
countable, so the product description gives the matching upper bound.
Part (iv) is Corollary~\ref{CD:countabletargets}; the final assertions
use Proposition~\ref{CD:sequential}.
\end{proof}

The criterion in (iv) is a condition on the support of the homotopy
groups, and makes no assertion that a general finite Postnikov complex
must be aspherical.  For a simply connected finite complex it reduces
to contractibility.  For example, Theorem~1.1 and its proof in
\cite{CDCasacuberta1998} make the reduced integral homology of such a
Postnikov piece rational; finite generation then forces it to vanish,
and the homological Whitehead theorem applies.  The unrestricted
fundamental-group issue is distinct, as the same primary account
explains.  No stronger asphericity claim is used here.

\subsection{Exactly which targets have a finite representative}

\begin{lemma}[Finite polyhedron comparison]
\label{CD:finitepoly}
Suppose $K$ is a finite category and $h_K$ inverts weak equivalences
between finite categories.  For every finite poset $P$, the canonical map
\[
             h_K(P)\longrightarrow[BP,BK]
\]
is bijective.
\end{lemma}
\begin{proof}
Put $A=NP$.  The simplicial set $A$ has finitely many nondegenerate
simplices.  Kan's replacement $\operatorname{Ex}^{\infty}NK$ is Kan
and weakly equivalent to $NK$; maps and homotopies from $A$ factor
through a finite iterate $\operatorname{Ex}^rNK$
\cite[\S\S4,6--8]{CDKan1957}.  By
$\operatorname{Sd}^r\dashv\operatorname{Ex}^r$, a map is therefore
represented by a functor $\operatorname{sd}^rP\to K$.
The last-vertex functor
$\lambda_r:\operatorname{sd}^rP\to P$ is a weak equivalence of finite
posets.  Invertibility of $h_K(\lambda_r)$ proves surjectivity.

For injectivity, a homotopy of two maps factors through some finite
iterate on $A\times\Delta[1]$.  Its adjoint is a functor
$H:\operatorname{sd}^r(P\times[1])\to K$ whose endpoints are the two
original functors precomposed with $\lambda_r$.
Let $e_i:\operatorname{sd}^rP\to
\operatorname{sd}^r(P\times[1])$ be the endpoint inclusions.
The square with last-vertex functors and
$\iota_i:P\to P\times[1]$ commutes.  In
$\mathrm{Cat}_{\mathrm{hot}}$, $\iota_0=\iota_1$ by the natural
transformation between them.  Since both last-vertex functors are
inverted by $h_K$, it follows that $h_K(e_0)=h_K(e_1)$.
Applying this equality to $H$ and then inverting $h_K(\lambda_r)$
proves injectivity.  This argument explicitly confines every category
used in the comparison to the finite test domain.
\end{proof}

\begin{lemma}[A compact-test equivalence with finitely generated source group]
\label{CD:anchoring}
Let $f:X\to Y$ induce bijections $[A,X]\to[A,Y]$ for all finite CW
complexes $A$.  If every component of $X$ has finitely generated
fundamental group, then $f$ is a weak equivalence.
\end{lemma}
\begin{proof}
A point identifies component sets, so work in corresponding connected
components and write $\phi:\Gamma\to G$ for the induced group map.
Injectivity on free circle classes implies $\ker\phi=1$.
For each $r$, the hypothesis on a wedge of $r$ circles is precisely
bijectivity of $\Gamma^r/\Gamma\to G^r/G$ for simultaneous conjugation.
Choose generators $\gamma_1,\ldots,\gamma_d$ of $\Gamma$.
For $g\in G$, surjectivity for $r=d+1$ supplies
$\gamma'_1,\ldots,\gamma'_d,a\in\Gamma$ and $h\in G$ with
\[
 \phi(\gamma_i)=h\phi(\gamma'_i)h^{-1},\qquad
 g=h\phi(a)h^{-1}.
\]
Injectivity for $r=d$ makes $(\gamma'_i)$ conjugate in $\Gamma$ to
$(\gamma_i)$.  Correcting $h$ by that conjugating element makes it
centralize $\phi(\Gamma)$, and the second equality then shows
$g\in\phi(\Gamma)$.  The case $d=0$ follows directly from circle
surjectivity.  Thus $\phi$ is an isomorphism.

For $j\ge2$, injectivity on free sphere classes makes the kernel of
$\pi_jf$ zero, since the orbit of zero is a singleton.  Surjectivity on
free sphere classes says that every element of $\pi_jY$ is a
$\pi_1Y$-translate of an element of the image.  That image is invariant
under the entire $\pi_1Y$ because $\pi_1f$ is onto.  Hence $\pi_jf$
is onto as well.  This proves the claim at every basepoint.
\end{proof}

\begin{lemma}[Polynomial exponents for elementary abelian tests]
\label{CD:elementarybound}
Let $E_r=(\Z/p)^r$.  If $T$ is a connected $h$-type with finite
homotopy groups and $h\ge1$, then
\[
 |[BE_r,T]|\ \le\
 |\pi_1T|^r
 \prod_{i=2}^{h}|\pi_iT|^{\binom{r+i-1}{i}}.
 \tag{\ref{CD:elementarybound}.1}
\]
For a nonzero finite abelian group $A$ with $p\mid |A|$ and $N\ge2$,
\[
                   |H^N(BE_r;A)|\ge p^{\binom rN}.
 \tag{\ref{CD:elementarybound}.2}
\]
\end{lemma}
\begin{proof}
Tensor the $r$ usual two-periodic free integral resolutions for $C_p$.
The resulting $\Z[E_r]$-resolution has free rank
$b_i=\binom{r+i-1}{i}$ in degree $i$, since weak compositions of $i$
into $r$ parts index its summands.  It is a resolution by the tensor
K\"unneth theorem: the factors are complexes of free abelian groups with
homology $\Z$ in degree zero only.  For every finite $E_r$-module $V$,
including nontrivial actions, this gives
$|H^i(E_r;V)|\le |V|^{b_i}$.
There are at most $|\pi_1T|^r$ fundamental-group homomorphisms.
With one chosen, successive based Postnikov lifts, when they exist,
form torsors under the appropriate local-coefficient $H^i$ before a
possible further quotient by homotopies.  Multiplying these bounds
proves (1); passing to free homotopy only decreases the count.

For (2), the integral K\"unneth splitting of
$(BC_p)^r$ shows that $H_N(BE_r;\Z)$ has a direct summand
$(\Z/p)^{\binom rN}$: use $H_1(BC_p;\Z)$ in $N$ distinct factors
and $H_0$ in the others.  The universal coefficient surjection
$H^N(BE_r;A)\twoheadrightarrow\Hom(H_N(BE_r;\Z),A)$,
and the existence of a nonzero element of $A[p]$, give the inequality.
\end{proof}

\begin{theorem}[Finite representatives: complete classification]
\label{CD:finiterepresentative}
For an arbitrary homotopy type $Y$, $F_Y$ is represented on finite
categories by a finite category if and only if
\[
                 Y\simeq\coprod_{c\in T}BG_c
\]
for a finite set $T$ and finite groups $G_c$.
The empty union is allowed.  Any finite representing category is
isomorphic in $\mathrm{Cat}_{\mathrm{hot}}$ to a finite groupoid.
For a finite CW target, a finite representative therefore exists
exactly when every component is contractible.
\end{theorem}
\begin{proof}
Suppose $\alpha:h_K\xrightarrow{\sim}F_Y$ with $K$ finite.
By Yoneda on $\mathcal C_f$, $\alpha$ is determined by
$y=\alpha_K(\id_K):BK\to Y$, and
$\alpha_C(u)=yBu$.  It inverts weak equivalences between finite tests.
Lemma~\ref{CD:finitepoly} shows that $y$ induces bijections on all finite
CW tests, because such a complex has a finite face-poset model.
The fundamental groups of $BK$ are finitely generated by its finite
one-skeleton.  Lemma~\ref{CD:anchoring} now proves $BK\simeq Y$.

Consequently $Y$ has finitely many components with finitely generated
fundamental groups.  Each set $[A,Y]$ for finite $A$ is finite, as it is
computed by a finite functor category.  If
$G=\pi_1Y_c$ is generated by $g_1,\ldots,g_d$, the classes of the tuples
$(g_1,\ldots,g_d,g)$, as $g$ ranges over $G$, are distinct under
simultaneous conjugation: a conjugator fixing the generating tuple
lies in $Z(G)$ and fixes $g$.  Finiteness of the test by $d+1$ circles
forces $G$ finite.  Finiteness of $[S^j,Y_c]$ then forces each
$\pi_jY_c$ finite.  Finally the six-element monoid and
Lemma~\ref{CD:wedgecriterion} force these groups to vanish for all
sufficiently large $j$.

It remains to rule out a top nonzero group $A=\pi_NY_c$ with $N\ge2$.
Choose $p\mid|A|$ and put $Z=P_{N-1}Y_c$.
Use a Postnikov fibration $K(A,N)\to Y_c\to Z$.
Over the constant map $BE_r\to Z$, the homotopy fibre of the mapping
space map is $\operatorname{Map}(BE_r,K(A,N))$.
Its components are $H^N(BE_r;A)$.  The only identifications of these
components in $[BE_r,Y_c]$ are orbits of
\[
 \pi_1\operatorname{Map}(BE_r,Z)_{\mathrm{const}}
       =[BE_r,\Omega Z].
\]
This is the path-lifting action in the homotopy exact sequence of a
fibration of mapping spaces.  Therefore
\[
 |[BE_r,Y_c]|\ \ge\
 \frac{|H^N(BE_r;A)|}{|[BE_r,\Omega Z]|}.
 \tag{\ref{CD:finiterepresentative}.1}
\]
All sets in the denominator are finite.  If $N=2$, $\Omega Z$ is a
finite discrete group.  If $N\ge3$, its finitely many components are
$(N-2)$-types with finite homotopy groups.  Applying
Lemma~\ref{CD:elementarybound} to them gives
\[
 \log_p|[BE_r,\Omega Z]|=O\!\left(r^{\max(1,N-2)}\right)
 \quad(N\ge3),
\]
and a constant bound for $N=2$.  The numerator in (1) is at least
$p^{\binom rN}$.  Thus its logarithm grows with positive leading term
of degree $N$, even after division by the denominator.
On the other hand a functor $\mathsf BE_r\to K$ chooses an object of
$K$ and the images of $r$ group generators.  With
$o=|\operatorname{Ob}K|$ and $m=|\Mor K|$,
\[
           |[BE_r,Y]|=|h_K(\mathsf BE_r)|\le o\,m^r.
\]
This has logarithm at most linear in $r$, a contradiction.
Every component of $Y$ is therefore $BG_c$ with $G_c$ finite.

Conversely choose a finite skeletal groupoid with exactly these groups.
Functors to a groupoid, and their natural transformations, extend
uniquely over groupoidification of the source.  Its functor-category
components are precisely free maps of classifying spaces, so it
represents $F_Y$ on all categories, in particular finite ones.
Both this groupoid and $K$ belong to $\mathcal C_f$; Yoneda makes the
two representatives isomorphic in $\mathrm{Cat}_{\mathrm{hot}}$.

Finally a nontrivial finite group cannot have a finite-dimensional
classifying space.  A subgroup $C_p$ would give a finite-dimensional
covering model for $BC_p$, contradicting the nonvanishing
$H^{2j}(C_p;\Z)=\Z/p$ for all $j\ge1$, computed by its periodic
resolution.  This proves the last assertion for finite CW targets.
\end{proof}

This proves the finite-target conclusions of R1 D.2 without a Schur
multiplier computation, and extends them to every $K(\Z/p,n)$ with
$n\ge2$.  Among the named connected targets in C1, the finite
representatives are exactly $BG$ for finite $G$ (including
$K(\Z/p,1)$).  The circle and every $K(\Z,n)$ have countable locally
finite representatives but no finite one.

\section{Finite stages with sphere colimit: the comparison can fail}
\label{sec:finitecomparison}

\begin{theorem}[Bounded homology suffices for countability]
\label{CD:boundedhomology}
Let $X$ be a connected CW complex whose integral homology groups are
finitely generated and vanish in sufficiently high degrees.
Then $[X,S^3]$ and $[X,S^2]$ are countable.
In particular, if $C$ is finite and $[BC,S^2]$ is uncountable,
$H_j(BC;\Z)$ is nonzero for arbitrarily large $j$.
\end{theorem}
\begin{proof}
The reduced suspension $\Sigma X$ is simply connected and has bounded
finitely generated integral homology.  Hatcher's Proposition~4C.1,
read with its proof on pp.~429--430 of the supplied text, constructs
a homotopy-equivalent CW complex with one generator cell per free
cyclic summand and one generator--relator pair per torsion cyclic
summand.  In this case there are only finitely many cells.  Thus
\[
 [X,S^3]=[X,\mathrm{SU}(2)]_*
       \cong[\Sigma X,B\mathrm{SU}(2)]_*
\]
is countable: $B\mathrm{SU}(2)\simeq\mathbb HP^\infty$ has a countable
CW model, and maps from a finite CW complex into such a model have
countably many classes.  Free and pointed classes agree for the
simply connected target $S^3$.

We verify the additional step for $S^2=\mathbb CP^1$, since an
injection from the Hopf lift set alone would not prove this direction.
For a map $f:X\to\mathbb CP^1$, let $L_f\subset X\times\mathbb C^2$
be its tautological complex line bundle.  Two maps $f_0,f_1$ with
the same $c_1\in H^2(X;\Z)$ have isomorphic line bundles.
Choose a unitary line-bundle isomorphism
$\varphi:L_{f_0}\to L_{f_1}$.
The fixed antilinear map
$J(z_1,z_2)=(-\overline z_2,\overline z_1)$ identifies any line with
its orthogonal complement.  Extend $\varphi$ to the complements by
$J\varphi J^{-1}$.  The resulting continuous unitary automorphism
of $X\times\mathbb C^2$ has determinant one, as is checked in the
frames $(z,Jz)$ and $(\varphi z,J\varphi z)$.
It is a map $g:X\to\mathrm{SU}(2)$ satisfying $f_1(x)=g(x)f_0(x)$.
Consequently $[X,S^3]$ acts transitively on each nonempty fibre of
\[
                  c_1:[X,S^2]\longrightarrow H^2(X;\Z).
\]
The latter group is countable by the universal coefficient theorem,
and the acting group is countable.  This proves countability.
For a finite category, homology is finitely generated in each degree
because its nerve has finitely many simplices in each degree.
Apply the connected assertion to its finitely many components.
\end{proof}

\begin{proposition}[Audit of the bounded-homology result R4]
\label{CD:monoidfinitehomology}
For a finite monoid $L$, bounded integral homology is equivalent to
$BL$ having finite simply connected CW type, and is equivalent to
$BL$ having finite-dimensional CW type.  If a prime $p$ divides its
universal group $U(L)$, then $H_j(BL;\mathbb F_p)$ is nonzero in
arbitrarily high degrees.
\end{proposition}
\begin{proof}
We check the proof of OWN\_RESULTS 1.1--1.2 and 2.1--2.2, including
the step involving an infinite-dimensional source.  For complex
$K$-theory put
\[
 F^q(X)=\ker\bigl(K^0(X)\to K^0(X^{(q-1)})\bigr).
\]
Relative products and a cellular approximation to the diagonal give
$F^aF^b\subseteq F^{a+b}$.  If $H^q(X;\Z)=0$, then
$F^q=F^{q+1}$.  Indeed the restriction of $x\in F^q$ to $X^{(q)}$
is represented by a relative cellular $K$-cochain $c$.  For odd $q$
the coefficient group is zero.  For even $q$, extension of $x$ over
$X^{(q+1)}$ says that the cellular coboundary of $c$ is zero; this is
the boundary in the exact sequence of the three successive skeleta.
Write $c=\delta b$ using $H^q(X;\Z)=0$.  The image of this relative
coboundary in $K^0(X^{(q)})$ is zero by the pair exact sequence.
Thus $x\in F^{q+1}$.  This is a skeletal argument and requires no
convergence assertion for an infinite spectral sequence.

Consequently if $H^q(X;\Z)=0$ for $q>n$, every element of
$F^{n+1}$ restricts to zero on every finite skeleton.  Suppose
$\pi_1X=\pi$ is finite, let $V$ be the flat bundle associated to
its regular representation $\rho$ of degree $d=|\pi|$, and put
$v=[V]-d\in F^1(X)$.  For $f:BG\to X$ with $G$ finite,
$f^*(v^{n+1})$ therefore vanishes on every skeleton of $BG$.
Atiyah's completion isomorphism and its kernel calculation
\cite[Propositions 6.10--6.11 and Theorem 7.2]{CDAtiyah1961}
say exactly that the character
$(\chi_{\rho f_*}-d)^{n+1}$ vanishes on every prime-power-order
element of $G$.  The regular character equals $d$ only at the
identity, so $f_*$ kills all such elements.  The primary components
of an element of a finite group generate its cyclic subgroup, hence
$f_*=1$.  Only the inverse limit of the skeletal $K$-groups is used;
no assertion that a phantom $K$-class is zero was made.

The prime-local variant used in OWN\_RESULTS 1.2 also checks.
Use the same skeletal argument for $KU_{(p)}$.  If $p$ divided the
order of $\operatorname{im}f_*$, lift an order-$p$ element and take
the $p$-primary part of its lift.  Restrict to the resulting cyclic
$p$-subgroup $C$.  Its odd lens-space skeleta have reduced $K^0$
a finite $p$-group, by the finite cellular Atiyah--Hirzebruch
filtration, so localization at $p$ is injective on their $K^0$.
The virtual character power is therefore zero in the integral
skeletal completion.  The cyclic proof of Atiyah's Proposition~8.1,
together with Proposition~6.11, makes the completion map on $R(C)$
injective.  Evaluating at the chosen element gives
$(-d)^{n+1}\ne0$, a contradiction.  Thus a bound on
$H^q(X;\Z_{(p)})$ excludes $p$ from the image order.

We also verify the semigroup input.  Every element of the finite
monoid $L$ has eventually periodic powers, so its image in $U(L)$
has an inverse represented by a positive power.  The image of $L$
is already a group and surjects onto $U(L)$; in particular $U(L)$
is finite.  Choose an idempotent $e$ in the minimal nonempty
two-sided ideal.  The finite monoid $eLe$ is a group: for any
$h\in eLe$, minimality gives $e=ahb$, and multiplying by $e$ at
both ends gives that equation inside $eLe$; in a finite monoid a
one-sided inverse is two-sided, which makes $h$ a unit.
Since $\eta(e)=1$, every $\eta(l)$ equals $\eta(ele)$, so
$eLe\to U(L)$ is onto.

There is a map $B(eLe)\to BL$ inducing this quotient.
For clarity, include $L$ in its idempotent completion and include
the group at the object $e$.  For each idempotent $e$, the comma
category of $L$ over $e$ has objects $a$ with $ea=a$ and arrows
$m:a\to b$ with $bm=a$.  The arrows $a:a\to e$ form a natural
transformation from its identity functor to the constant functor
at $e$; thus it is contractible.  Quillen's Theorem~A makes the
inclusion of $L$ a weak equivalence and supplies the claimed map.
The preceding regular-representation argument now says that if
the integral cohomology of $BL$ is bounded, this surjective
fundamental-group map is trivial.  Hence $U(L)=1$.

The normalized bar chain groups of a finite monoid are finitely
generated in each degree.  Bounded integral homology therefore
gives bounded cohomology by the universal coefficient theorem.
We have just proved simple connectivity; Hatcher's verified
Proposition~4C.1 gives a finite CW model.  The reverse implications
are immediate from cellular homology.  If mod-$p$ homology were
bounded, finite generation and the universal coefficient theorem
would instead give bounded cohomology with $\Z_{(p)}$ coefficients.
The prime-local argument would imply $p\nmid|U(L)|$, proving the
last assertion.  This also checks OWN\_RESULTS 2.3 and 2.4: a
finite-length projective resolution forces the same Tor vanishing.
\end{proof}

The finitely dominated assertion in R4 needs less: if $X$ is a
homotopy retract of a finite CW complex $A$, then $[X,S^2]$ injects
into $[A,S^2]$ by precomposition with the retraction.  The latter set
is countable.  This verifies that assertion, without assuming a
finitely dominated space has finite CW type.

\begin{lemma}[A finite acyclic approximation of every finite category]
\label{CD:spread}
For a nonempty finite category $C$, define $P_n(C)$ by
\[
 \operatorname{Ob}P_n(C)=\{(i,c):0\le i\le n,\ c\in\operatorname{Ob}C\}.
\]
There is only the identity at each object at a fixed level; for $i<j$,
the arrows $(i,c)\to(j,d)$ are $C(c,d)$, with composition induced by
$C$.  Let $p_n:P_n(C)\to C$ forget the level and retain the arrow label.
Then $P_n(C)$ is finite and acyclic, $BP_n(C)$ has dimension at most $n$,
and the induced functor
\[
          p_\infty:P_\infty(C)=\bigcup_nP_n(C)\longrightarrow C
\]
is a weak equivalence.
\end{lemma}
\begin{proof}
Strict increase of levels proves the dimension bound.
For $c\in C$, an object of $(p_\infty\downarrow c)$ is
$(i,d,a:d\to c)$.  Finitely many such objects admit arrows to
$(L,c,\id_c)$ for a larger level $L$, using their labels $a$.
Parallel arrows into $(j,d,b)$ become equal after composition with
the arrow labelled $b$ to $(L,c,\id_c)$.
Thus this comma category is nonempty and filtered, so its nerve is
contractible.  Quillen's Theorem~A applies to the specified
$p_\infty$ and proves the claimed equivalence.
\end{proof}

\begin{theorem}[C3, with finite categories at every stage]
\label{CD:comparisonfailure}
Let $C$ be any finite category for which $H_j(BC;\Z)$ is nonzero
in arbitrarily large degrees.  There is an explicitly defined sequence
of finite categories $J_n$ with
$B\colim_nJ_n\simeq S^2$ for which the canonical map
\[
 \colim_n[BC,BJ_n]\longrightarrow[BC,B\colim_nJ_n]
 \tag{\ref{CD:comparisonfailure}.1}
\]
is not injective.  Therefore such a sequence exists for every finite
$C$ with uncountable $[BC,S^2]$.
\end{theorem}
\begin{proof}
Let $K_n$ be the Grothendieck construction of the span
\[
                         *\longleftarrow P_n(C)\xrightarrow{p_n}C.
\]
All categories and all cross-arrow sets are finite.  The inclusions
$P_n(C)\to P_{n+1}(C)$ give actual transition functors $K_n\to K_{n+1}$
that are the identity on the copies of $C$ and $*$.
The categorical mapping-cylinder comparison gives a natural weak
equivalence
\[
                   BK_n\simeq\Cof(Bp_n:BP_n(C)\to BC),
\]
compatible with the map $a_n:BC\to BK_n$ from the copy of $C$.
One can check this comparison without a general homotopy-colimit
theorem.  A half of the span construction has objects $A\amalg B$
and arrows from $a\in A$ to $b\in B$ equal to $B(Fa,b)$.
Its projection to $B$ is left adjoint to the inclusion of $B$, so
its realization retracts up to homotopy to $BB$, with the prescribed
restriction $BF$ on $BA$.  The two halves have intersection $BA$,
a simplicial subcomplex.  Their union consequently computes the
homotopy pushout of the span.  This also proves naturality here.

The relative homology sequence and $\dim BP_n(C)\le n$ show that
\[
 (a_n)_*:H_k(BC;\Z)\xrightarrow{\sim}H_k(BK_n;\Z)
                         \qquad(k>n+1).
 \tag{\ref{CD:comparisonfailure}.2}
\]
The same span description for $K_\infty=\colim_nK_n$ and
Lemma~\ref{CD:spread} give
$BK_\infty\simeq\Cof(Bp_\infty)\simeq *$.
Nerves commute with these filtered colimits: a simplex is a finite
composable string, and all its equalities occur at one stage.

Fix a finite category $Q$ of type $S^2$, for example the five-element
rectangular-band monoid with identity, and fix objects $c_0\in C$,
$q_0\in Q$.  Let $J_n$ be the Grothendieck construction of the span
\[
                 K_n\longleftarrow *\longrightarrow Q,
\]
whose maps choose the copy of $c_0$ and $q_0$.
The transition functors are again actual functors of finite categories.
The preceding mapping-cylinder description identifies
$BJ_n$ with the wedge $BK_n\vee BQ$ up to a natural weak equivalence.
There is an actual retraction $J_n\to Q$, constant on $K_n$, and its
colimit realizes a weak equivalence
$BJ_\infty\to BQ\simeq S^2$ because $BK_\infty$ is contractible.

Compare the functor $\alpha_0:C\to K_0\to J_0$ with the constant
functor $\beta_0:C\to Q\to J_0$ at $q_0$.
At each stage their realizations remain different: choose
$k>\max(n+1,2)$ with $H_k(BC;\Z)\ne0$ and use (2) and the wedge
splitting.  The map $B\alpha_n$ is nonzero on that homology group,
whereas $B\beta_n$ is zero.
In the colimit $B\alpha_\infty$ factors through the contractible
$BK_\infty$, and it is homotopic to $B\beta_\infty$.
Thus the two classes are distinct in the sequential colimit on the
left of (1) and equal on the right, as asserted.
The final sentence follows from Theorem~\ref{CD:boundedhomology}.
\end{proof}

For the six-element witness, the program
\texttt{code/CD\_finite\_telescope.py} constructs the actual multiplication
tables of $P_n,K_n,J_n$ for $0\le n\le5$.  It checks all identity laws,
all composable associativity triples, every transition functor, and
the two comparison maps and retractions.  The $J_n$ in these checks
have respectively $32,46,66,92,124,162$ arrows.  The all-degree
nontriviality and the colimit equivalence come from the proof above,
not from extrapolation of these finite checks.

\begin{proposition}[What a sphere colimit forces]
\label{CD:surjectivity}
For any actual filtered diagram of categories whose colimit has finite
CW type, the comparison
\[
             \colim_i[X,BK_i]\longrightarrow[X,B\colim_iK_i]
\]
is onto for every CW complex $X$.  If every $BK_i$ also has finite
CW type, the comparison is bijective for every $X$.
In particular, in C3 with sphere colimit the failure must be one of
injectivity, and a sequence consisting only of finite posets cannot
supply that failure.
\end{proposition}
\begin{proof}
Put $Y=B\colim_iK_i\simeq\hocolim_iBK_i$.
Compactness of a finite CW model of $Y$ gives a map
$s:Y\to BK_j$ with $\iota_js\simeq\id_Y$.
Thus every $f:X\to Y$ is represented at stage $j$ by $sf$.
For injectivity, put two representatives at one common stage $k$.
After taking a common successor of $j,k$, compare the transition
$BK_k\to BK_\ell$ with $BK_k\xrightarrow{\iota_k}Y
\xrightarrow{s}BK_j\to BK_\ell$.
Their composites to $Y$ agree.  If $BK_k$ has finite CW type,
compactness makes these two maps equal at a further stage.
Two maps from $X$ whose composites to $Y$ agree therefore become
equal there as well.  Finite posets have finite CW nerves.
\end{proof}

\begin{proposition}[Direct audit of R1 E.8]
\label{CD:E8}
For an actual filtered diagram $D:I\to\mathrm{Cat}$, the canonical
functor $\int_I D\to\colim_I D$ is a weak equivalence.
Objectwise weak equivalences of such diagrams induce weak
equivalences of colimits.  The comparison on mapping classes out
of a finite CW complex is bijective, but the assertion with an
arbitrary finite-category nerve as source is false, as in R1 E.8.
\end{proposition}
\begin{proof}
Here is a direct verification of the specific first comparison.
For an object $x$ of the colimit, the comma category
$(\int_I D\downarrow x)$ is filtered.  Any finite list of its
objects consists of objects at stages and arrows from their colimit
images to $x$.  Lift these arrows and a representative of $x$ to a
common stage.  They give arrows to the comma object represented by
that representative and its identity.  For two parallel comma
arrows, first equalize their finitely many index maps in $I$;
their equality after composition with the structure arrow to $x$
is then witnessed at a later stage.  That structure arrow supplies
the required equalizer.  Nonemptiness follows by representing $x$.
Its nerve is contractible, so Quillen's Theorem~A proves the claim.

The nerve commutes with filtered colimits because each $[n]$ is a
finitely presented category.  Kan's $\operatorname{Ex}^{\infty}$
also commutes with them: each subdivided standard simplex is finite.
A filtered colimit of Kan complexes is Kan, and its homotopy groups
are the filtered colimits of their homotopy groups, since spheres,
horns and homotopies involve finite simplicial data.  This proves
the assertion about objectwise weak equivalences.  The same finite
data argument on a finite simplicial model of a finite CW complex
proves the mapping-class comparison, for both maps and homotopies.
This checks the compactness step of E.8 without extending it to an
infinite-dimensional nerve.

For its stated counterexample put
$D_r=c\operatorname{Sd}^2\operatorname{sk}_rN(\mathsf BC_2)$.
These are finite cofibrant posets: finiteness of the vertex set and
Thomason's poset lemma give finiteness of all arrows
\cite[Proposition 4.6 and Lemma 5.6]{CDThomason1980}.
The canonical subdivision comparisons are natural weak equivalences
\cite[Theorem 4.11 and Example 4.12(iv)--(v)]{CDFritschLatch1981},
so $BD_r\simeq|\operatorname{sk}_rN(\mathsf BC_2)|$ and
$B\colim_rD_r\simeq BC_2$.  The transitions are actual functors,
as required.  The identity of $BC_2$ does not factor, even up to
homotopy, through any $BD_r$: such a factorization would make
$H^*(BC_2;\mathbb F_2)$ a retract of the cohomology of a finite
CW complex.  The first group is nonzero in every degree, whereas
the latter vanishes above its dimension.  Thus this comparison
fails surjectivity.  In contrast, the sphere-colimit counterexamples
of Theorem~\ref{CD:comparisonfailure} fail injectivity, as
Proposition~\ref{CD:surjectivity} requires.
\end{proof}

\section{Consequences at the specified manuscript pages}
\label{sec:manuscript-CD}

\begin{theorem}[D: the corrected finite-test statements]
\label{CD:manuscript}
At V~ms.~43--46, V~ms.~51--52 and VII~ms.~121--124, 129--132,
the following replacements are valid.
\begin{enumerate}[label=(\roman*)]
\item On finite categories, every $[B-,Y]$ admits a filtered
presentation by finite-category representables.  A sequence exists
exactly in Theorem~\ref{CD:alltargets}; for countable CW targets this
means componentwise Postnikov truncation.  In particular there is no
such sequence for $S^2$ or for any $S^n$, $n\ge2$, or $\mathbb CP^2$.
\item For an actual sequence of finite categories, its ordinary
categorical colimit is countable and locally finite, and
$\colim_n h_{K_n}(C)=h_{\colim_nK_n}(C)$ on finite $C$.
This is the valid categorical finite-presentation assertion at
V~ms.~45.  It does not replace $h_K(C)$ by $[BC,BK]$.
\item A finite representative exists exactly for the finite groupoid
types of Theorem~\ref{CD:finiterepresentative}.  Any sequential
presentation of another representable target functor that does exist
must have unbounded stage sizes, by Proposition~\ref{CD:uniformfinite}.
\item The claim at VII~ms.~121 that homotopy classes out of every
finite category commute with filtered colimits is false.  Compact
CW tests give the usual valid statement.  Even with finite stages
and sphere colimit, Theorem~\ref{CD:comparisonfailure} gives a failure
for every finite $C$ with uncountable $[BC,S^2]$.
\item Set-valued tests by finite categories do not form a conservative
restricted Yoneda functor on all homotopy types.  Consequently the
conservativity step at V~ms.~52 and VII~ms.~121, and conclusions
at VII~ms.~123--124 relying on that step, cannot be used.
\item The invariance condition at VII~ms.~129--132 must specify its
test domain.  Restricted to finite categories it permits all the
functors in (i), including higher Eilenberg--Mac Lane targets.  It does
not impose the global groupoid conclusion merely by referring to
finite tests.  The exact finite-domain extension conditions in
R5 B.8 are verified above.
\end{enumerate}
\end{theorem}
\begin{proof}
Statements (i)--(iv) follow from the results explicitly cited in them.
For (v) a concrete counterexample uses
$G=A_{\mathrm{fin}}(\mathbb N)$ and the shift embedding
$\sigma:G\to G$ whose image fixes $0$.
It is not onto.  Nevertheless $B\sigma$ induces a bijection on
$[BC,BG]$ for every finite category $C$: the finite-support argument
in Remark~\ref{CD:largegroups} applies to
$\mathbb N\setminus\{0\}\subset\mathbb N$.
The map is therefore an isomorphism on all the proposed set-valued
finite tests and is not a weak equivalence, since its fundamental-group
map is not surjective.  This supplies an explicit failure of
conservativity with countable source and target.
Finally $K(\Z/p,2)$ has finite values on finite tests, by the
cochain calculation, and so has the presentation in (i), although it
is not a groupoid type and has no finite representative.  Thus a
global groupoid conclusion cannot follow from the finite-domain
invariance condition in (vi).
\end{proof}

The manuscript formulas at V~ms.~45 and VII~ms.~131 involving
constant diagrams should have matching source and target subscripts:
$\operatorname{Hom}(a_X,b_X)$ at V~ms.~45 and
$\operatorname{Hom}(a_Y,b_Y)$ when the test category is named $Y$.
Those transcription corrections do not repair the separate
compactness and conservativity errors proved above.
